\documentclass[11pt,a4paper]{article}
\usepackage{harmonikabau_preamble}

% == Document Metadata ====================
\newcommand{\hbdoknr}{0017}
\newcommand{\hbtitel}{Treble Reeds}
\newcommand{\hbuntertitel}{Overtone Modes F3 to C6}
\newcommand{\hbheadtext}{Doc.\,0017 --- Treble Reeds: Overtone Modes F3--C6}
\newcommand{\hbfoottext}{Cross-references: Doc.\,0016}

\AtBeginDocument{%
  \fancyhead[L]{\small\hbheadtext}%
  \fancyhead[R]{\small\thepage}%
  \fancyfoot[C]{\small\color{hb@darkgray}\hbfoottext}%
}

\begin{document}
\hbmaketitle
\begin{tcolorbox}[colframe=red!60!black, colback=red!5!white]
\textbf{Note:} This document presents an explanatory model.
The interpolated scale is an estimate --- real reed plates
vary depending on the manufacturer.
\end{tcolorbox}

\noindent\textit{\small Calculation script: The complete FEM calculation with all tables
is deposited in the same directory on GitHub:
\textbf{pruefskript\_0017\_diskant.py} (W independence, uniform minimum thickness,
calibrated profiling, $f_2$ vs.\ $f_H$, bass/treble comparison).}

\noindent\textit{FEM calculation of the bending modes of treble reeds with variable
scale (L = 37$\to$14\,mm, W = 3.7$\to$1.7\,mm). Calibrated thickness for each
semitone. Comparison with the bass reeds from Doc.~0016.}

\section{Width W Has No Effect on Pitch}

For an Euler--Bernoulli cantilever:
\begin{equation}
f = \frac{\beta^2}{2\pi L^2} \times t \times \sqrt{\frac{E}{12\rho}}
\end{equation}

Width $W$ cancels out: $I = Wt^3/12$, $A = Wt$, so $I/A = t^2/12$.
Pitch is determined solely by $t$ and $L$.
$W$ influences mass ($\propto W$), volume flow ($\propto W$),
loudness and torsional stiffness ($\propto W^3$) --- but not the pitch.

\begin{tcolorbox}[colframe=keygreen, colback=green!5!white]
\textbf{W determines loudness and response, not pitch.}
Hence $W$ and $L$ decrease steadily from low to high ---
the low notes require more volume flow.
\end{tcolorbox}

\section{Treble Scale: L and W Decrease Steadily, t is Profiled}

Unlike the bass, in the treble both length and width vary continuously.
The foot thickness ($t_0$) is always slightly greater than the theoretical
minimum thickness --- profiling (thinning towards the tip) adjusts the pitch.

\begin{center}
\resizebox{\linewidth}{!}{\begin{tabular}{lcccccccc}
\hbtoprule
Note & $f$ & $L$ & $W$ & $t_0$ & $t_\mathrm{tip}$ & Thin. & $f_2/f_1$ & $f_2$ \\
\midrule
F3 & 175 & 37.0 & 3.70 & 0.298 & 0.223 & 28\,\% & 5.32 & 348 \\
C3 & 131 & 29.0 & 3.05 & 0.133 & 0.104 & 24\,\% & 5.33 & 697 \\
C4 & 262 & 22.8 & 2.51 & 0.161 & 0.132 & 20\,\% & 5.45 & 1426 \\
C5 & 523 & 17.9 & 2.06 & 0.197 & 0.169 & 16\,\% & 5.63 & 2947 \\
C6 & 1047 & 14.0 & 1.70 & 0.244 & 0.217 & 12\,\% & 5.82 & 6094 \\
\bottomrule
\end{tabular}}
\end{center}

\begin{tcolorbox}[colframe=keygreen, colback=green!5!white]
\textbf{$f_2/f_1 = 5{.}3$--$5{.}8$} (profiled) --- lower than the 6.27
of the uniform reed, because profiling compresses the modes.
The light profiling in the treble (12--28\,\%) lowers the ratio
less than in the bass (30--50\,\%).
\end{tcolorbox}

\section{Absolute $f_2$: Where Does Mode~2 Hit the Chamber Resonance?}

In the bass the critical zone is A1--A2 ($f_2 = 345$--$689$\,Hz,
$f_H$ bass $= 200$--$700$\,Hz). In the treble everything shifts upward:
$f_2$ lies for C2--C3 at 410--820\,Hz --- \textbf{below} the treble chambers.
Only from C4 ($f_2 \approx 1640$\,Hz) does $f_2$ begin to hit the treble chamber range.
The critical zone in the treble is \textbf{C4--C5}
($f_2 = 1640$--$3280$\,Hz).

\section{Profiling in the Treble: Less Room for Adjustment}

Treble reeds are lightly profiled: 12--25\,\% thinning
(bass reeds: 30--50\,\%). The thinnest point is at the tip for the higher notes,
and somewhat towards the middle for the lower ones. Profiling reduces
$f_2/f_1$ by 0.4--0.9 (from 6.27 to 5.4--5.8). In the bass the reduction was
1.0--2.3. Less room for adjustment with thinner reeds.

\section{Stiffness and Mass}

Stiffness increases exponentially: factor~100 from C2 (4.8\,N/m) to
C6 (489\,N/m). Mass decreases from 117\,mg to 47\,mg. High notes:
stiffer, lighter, higher threshold pressure, faster onset.

\section{Comparison: Bass vs.\ Treble}

\textbf{Common:} $f_2/f_1 = 6{.}27$ = constant. The cantilever physics
are identical.

\textbf{Difference 1 --- Scale:} In the bass, one reed length covers 2~octaves.
In the treble, $L$, $W$ and $t$ all vary simultaneously.

\textbf{Difference 2 --- Profiling:} Bass: 30--50\,\%, effect $-1{.}0$ to
$-2{.}3$. Treble: 12--25\,\%, effect $-0{.}4$ to $-0{.}9$.

\textbf{Difference 3 --- Critical zone:} Bass: A1--A2 ($f_2$ in the $f_H$ range
of the bass chambers). Treble: C4--C5 ($f_2$ in the $f_H$ range of the treble chambers).
The lower treble notes (F3--B3) are non-critical.

\section{Summary}

\begin{enumerate}
\item \textbf{W has no effect on pitch.} $f \propto t/L^2$.
  $W$ determines loudness and response.
\item \textbf{$f_2/f_1 = 6{.}27$ (uniform) = identical to bass.}
  Profiled: 5.3--5.8 depending on thinning depth.
\item \textbf{Actual foot thickness $>$ minimum thickness.} Profiling
  raises $f_1$ --- the tuner starts thicker and grinds down.
\item \textbf{Critical zone: C4--C5.} $f_2 = 1400$--$3000$\,Hz hits
  the treble chambers. Lower notes (F3--B3) non-critical.
\item \textbf{Stiffness:} Factor~100 over 4~octaves.
\end{enumerate}

\bigskip
\textit{Bass and treble obey the same physics.
The scale changes the absolute frequencies --- the ratios remain.}


\end{document}
